\documentclass[border=3pt,tikz]{standalone}
\usepackage{amsmath} % for \text
\usepackage{etoolbox} % ifthen
\usepackage[outline]{contour} % glow around text
\usetikzlibrary{calc} % for adding up coordinates
\usetikzlibrary{decorations.markings,decorations.pathmorphing}
\usetikzlibrary{angles,quotes} % for pic (angle labels)
\usetikzlibrary{arrows.meta} % for arrow size
\usepackage{xfp} % higher precision (16 digits?)
\contourlength{1.1pt}

% COLORS
\colorlet{myred}{red!85!black}
\colorlet{mydarkred}{red!55!black}
\colorlet{mylightred}{red!85!black!12}
\colorlet{myfieldred}{mydarkred!5} % for S' background
\colorlet{myredhighlight}{myred!20} % highlights simultaneity in ladder paradox
\colorlet{myblue}{blue!80!black}
\colorlet{mydarkblue}{blue!50!black}
\colorlet{mylightblue}{blue!50!black!30}
\colorlet{mylightblue2}{myblue!10}
\colorlet{mygreen}{green!80!black}
\colorlet{mydarkgreen}{green!50!black}
\colorlet{mypurple}{blue!40!red!80!black}
\colorlet{mydarkpurple}{blue!40!red!50!black}
\colorlet{mypurple2}{blue!70!red!90!black}
\colorlet{mydarkpurple2}{blue!70!red!60!black}
\colorlet{myorange}{orange!40!yellow!95!black}
\colorlet{mydarkorange}{orange!40!yellow!85!black}
\colorlet{mybrown}{brown!20!orange!90!black}
\colorlet{mydarkbrown}{brown!20!orange!55!black}
\colorlet{mypurplehighlight}{mydarkpurple!20} % highlights simultaneity in ladder paradox

% STYLES
\tikzset{
	>=latex, % for LaTeX arrow head
	world line/.style={myblue!40,line width=0.2},
	world line t/.style={mypurple!50!myblue!40,line width=0.2},
	world line'/.style={mydarkred!40,line width=0.2},
	mysmallarr/.style={-{Latex[length=3,width=2]},thin},
	mydashed/.style={dash pattern=on 3 off 3},
	rod/.style={mydarkbrown,draw=mydarkbrown,double=mybrown,double distance=2pt,
		line width=0.2,line cap=round,shorten >=1pt,shorten <=1pt},
	%rod'/.style={rod,draw=mydarkbrown!80!red!85,double=mybrown!80!red!85},
	vector/.style={->,line width=1,line cap=round},
	vector'/.style={vector,shorten >=1.2},
	particle/.style={mygreen,line width=0.9},
	photon/.style={-{Latex[length=5,width=4]},myorange,line width=0.8,decorate,
		decoration={snake,amplitude=1.0,segment length=5,post length=5}}
}
\def\tick#1#2{\draw[thick] (#1) ++ (#2:0.06) --++ (#2-180:0.12)}
\def\tickp#1#2{\draw[thick,mydarkred] (#1) ++ (#2:0.06) --++ (#2-180:0.12)}
\def\Nsamples{100} % number samples in plot
\begin{document}
	
			
			
			% SPACETIME DIAGRAM - INVARIANT HYPERBOLOIDS with equations
			\def\drawaxes{ % common axes
				
				% SETTINGS
				\def\xmin{0.3}
				\def\xmax{3.1}
				\def\ymin{0.3}
				\def\ymax{2.6}
				\def\xminp{0.4} % minimum of rotated axis
				\def\xmaxp{2.85} % maximum of rotated axis
				\def\Nlines{4} % number of world lines (at constant x/t)
				\pgfmathsetmacro\ang{atan(0.5)} % angle between x and x' axes
				\pgfmathsetmacro\d{0.64*\xmax/\Nlines} % grid size
				\pgfmathsetmacro\D{\d/cos(\ang)/sqrt(1-tan(\ang)^2)} % grid size, boosted
				\pgfmathsetmacro\dextra{(\Nlines+1)*\d} % extra line
				\pgfmathsetmacro\st{3*\d} % spacetime interval
				\pgfmathsetmacro\sx{4*\d} % spacetime interval
				\pgfmathsetmacro\sr{sqrt(\sx^2-\st^2)} % spacetime interval sr^2 = st^2 - sx^2 < 0
				\pgfmathsetmacro\Ax{3*\D*sin(\ang)} % x coordinate of event A
				\pgfmathsetmacro\Ay{3*\D*cos(\ang)} % y coordinate of event A
				\pgfmathsetmacro\Bx{4*\D*cos(\ang)} % x coordinate of event B
				\pgfmathsetmacro\By{4*\D*sin(\ang)} % y coordinate of event B
				\pgfmathsetmacro\Cx{\Ay+\By} % x coordinate of event C' %(\Bx+tan(\ang)*\Ay)/sqrt(1-tan(\ang)^2)
				
				% COORDINATES
				\coordinate (O)  at (0,0);
				\coordinate (X)  at (\xmax+0.2,0);
				\coordinate (T)  at (0,\ymax+0.2);
				\coordinate (X') at (\ang:\xmaxp+0.2);
				\coordinate (T') at (90-\ang:\xmaxp+0.2);
				\coordinate (A)  at (0,\st);        % event A
				\coordinate (A') at (90-\ang:3*\D); % event A', boosted A
				\coordinate (B)  at (\sx,0);        % event A
				\coordinate (B') at (\ang:4*\D);    % event A', boosted A
				\coordinate (C)  at (4*\d,3*\d);    % event C
				
				% WORLD LINE GRID
				\message{  Making world lines...^^J}
				\foreach \i [evaluate={\x=\i*\d;}] in {1,...,\Nlines}{
					\message{  Running i/N=\i/\Nlines, x=\x...^^J}
					\draw[world line]   ( \x,-\ymin) -- ( \x,\ymax);
					\draw[world line t] (-\xmin, \x) -- (\xmax, \x);
				}
				\draw[world line]  (\dextra,-\ymin)  -- (\dextra,\ymax);
				\draw[world line]  (\dextra+\d,-\ymin) -- (\dextra+\d,\ymax);
				\draw[world line'] (-\xmin,\dextra) -- (\xmax,\dextra);
				
				% BOOSTED WORLD LINE GRID
				\message{  Making world lines for boosted frame...^^J}
				\fill[mydarkred,opacity=0.05]
				(O) --++ (\ang:\xmaxp) --++ (90-\ang:\xmaxp) --++ (\ang:-\xmaxp) -- cycle;
				\fill[mydarkred,opacity=0.05]
				(O) --++ (\ang-180:\xminp) --++ (-90-\ang:\xminp) --++ (\ang:\xminp) -- cycle;
				\foreach \i [evaluate={\x=\i*\D;}] in {1,...,\Nlines}{
					\message{  Running i/N=\i/\Nlines, x=\x...^^J}
					\draw[world line'] (\ang:\x) --++ (90-\ang:\xmaxp);
					\draw[world line'] (90-\ang:\x) --++ (\ang:\xmaxp);
				}
				
				% AXES
				\draw[->,thick] (0,-\ymin) -- (T) node[left=-1] {$ct$};
				\draw[->,thick] (-\xmin,0) -- (X) node[below=0] {$x$};
				\draw[->,thick,mydarkred] (-90-\ang:\xminp) -- (T')
				node[right=2,above=-1] {$ct'$};
				\draw[->,thick,mydarkred] (\ang-180:\xminp) -- (X') node[below=2,right=-3] {$x'$};
				
				% LIGHTCONE
				\draw[myorange,thick]
				(-1.1*\xminp,1.1*\xminp) -- (1.1*\xminp,-1.1*\xminp)
				(-1.3*\xminp,-1.3*\xminp) -- (1.2*\xmaxp,1.2*\xmaxp);
				
				% SPACELIKE HYPERBOLOIDS
				\draw[mygreen,thick,samples=\Nsamples,smooth,variable=\x,domain=-1.6*\xmin:1.05*\xmax]
				plot(\x,{sqrt((\st)^2+(\x)^2)});
				\draw[mydarkgreen,very thick,samples=\Nsamples,variable=\x,domain=0:\Ax,
				decoration={markings,mark=at position 0.58 with {\arrow{latex}}},postaction={decorate}]
				plot(\x,{sqrt((\st)^2+(\x)^2)});
				\node[mydarkgreen,above left=-2] at (\xmax,{sqrt((\st)^2+(\xmax)^2)})
				{$s^2 = c^2t^2-x^2>0$};
				
				% TIMELIKE HYPERBOLOIDS
				\draw[myred,very thick,samples=\Nsamples,variable=\y,domain=\st:\Cx,
				decoration={markings,mark=at position 0.58 with {\arrow{latex}}},postaction={decorate}]
				plot({sqrt(\sr^2+(\y)^2)},\y);
				\draw[myblue,thick,samples=\Nsamples,smooth,variable=\y,domain=-1.6*\ymin:0.95*\ymax]
				plot({sqrt(\sx^2+(\y)^2)},\y);
				\draw[mydarkblue,very thick,samples=\Nsamples,variable=\y,domain=0:\By,
				decoration={markings,mark=at position 0.58 with {\arrow{latex}}},postaction={decorate}]
				plot({sqrt(\sx^2+(\y)^2)},\y);
				\node[mydarkblue,below=0] at (0.7*\xmax,-1.6*\xmin)
				{$s^2 = c^2t^2-x^2<0$};
				
				% TICKS
				\draw[mydarkgreen,dashed] ({\Ax},0) -- (A') -- (0,{\Ay});
				\draw[mydarkblue,dashed] ({\Bx},0) -- (B') -- (0,{\By});
				\tick{0,\st}{0} node[mydarkgreen,right=4,below left=-2] {$ct_1$};
				\tick{\sx,0}{90} node[mydarkblue,below=1,below left=-3] {$x_1$};
				
				% EVENTS
				\fill[mydarkgreen] (A)  circle(0.03); % event A
				\fill[mydarkgreen] (A') circle(0.03); % event A'
				\fill[mydarkblue]  (B)  circle(0.03); % event B
				\fill[mydarkblue]  (B') circle(0.03); % event B'
				\fill[mydarkred]   (C)  circle(0.03); % event C
				\fill[mydarkred] (\ang:4*\D)++(90-\ang:3*\D) coordinate (C') circle(0.03); % event C'
			}
			\begin{tikzpicture}[scale=2]
				\message{Invariant hyperboloids with equations^^J}
				
				% AXES
				\drawaxes
				\node[mydarkgreen,below=1] at (\Ax/2,{sqrt((\st)^2+(\Ax/2)^2)}) {$\phi$};
				\node[mydarkblue,left=1.5] at ({sqrt(\sx^2+(0.54*\By)^2)},0.54*\By) {\contour{white}{$\phi$}};
				
				% TICKS
				\tick{0,\Ay}{0} node[mydarkgreen,above=0,left=-2]
				{$ct_1\cosh\phi$};
				\tick{\Ax,0}{90} node[mydarkgreen,right=4,below=-4]
				{\contour{white}{$ct_1\sinh\phi$}};
				\tick{\Bx,0}{90} node[mydarkblue,right=8,below=-4]
				{\contour{white}{$x_1\cosh\phi$}};
				\tick{0,\By}{0} node[mydarkblue,below=0,left=-2]
				{$x_1\sinh\phi$};
				
				% EVENT LABELS
				\node[mydarkred,anchor=0,inner sep=3] at (C) {\contour{myfieldred}{R}};
				\node[mydarkred,below right] at (1.0*\xmax,3.57*\d) {$
					\begin{aligned}
						%\mathrm{R} &= (x_1,ct_1) \\
						%           &= (x_1\cosh\phi-ct_1\sinh\phi,\\[-0.3em]
						%           &\hspace{1.7em} ct_1\cosh\phi-x_1\sinh\phi)\\
						\mathrm{R}
						&=
						\left\{\begin{aligned}
							ct &= ct_1 \\
							x &=  x_1
						\end{aligned}\right.\\
						&=
						\left\{\begin{aligned}
							ct' &= ct_1\cosh\phi -  x_1\sinh\phi \\
							x' &=  x_1\cosh\phi - ct_1\sinh\phi
						\end{aligned}\right.
					\end{aligned}
					$};
				%\node[mydarkred,anchor=-173,inner sep=3] at (C') {\contour{myfieldred}{R$'$}};
				\node[mydarkred,anchor=167,inner sep=3] at (C') {$
					\begin{aligned}
						\mathrm{\contour{myfieldred}{R$'$}}
						&=
						\left\{\begin{aligned}
							ct &= ct_1\cosh\phi +  x_1\sinh\phi \\
							x &=  x_1\cosh\phi + ct_1\sinh\phi
						\end{aligned}\right.\\
						&=
						\left\{\begin{aligned}
							ct' &= ct_1 \\
							x' &=  x_1
						\end{aligned}\right.
					\end{aligned}
					$};
				
			\end{tikzpicture}
		\end{document}